\chapter{相关工作}

\section{移动边缘计算与无人机辅助计算}
MEC 通过将计算能力部署在接入侧，缩短终端任务到远端云的传输路径，其资源优化通常同时涉及无线传输、边缘计算、任务准入和能耗约束~\cite{mao2017survey}。无人机作为可快速部署的空中接入与计算节点，又引入覆盖、高度、轨迹和飞行能耗等变量~\cite{zeng2016uav,wu2018joint}。近年的多 UAV-MEC 研究多把轨迹、计算卸载和能量补给纳入长期联合决策，例如采用多智能体深度强化学习联合优化任务卸载与航迹~\cite{ju2024multi}，或联合考虑卸载、无线充电与 UAV 协作调度~\cite{zhao2024collaborative}。这类工作回答的是动态部署或长期运行问题，而本文主动固定 UAV 悬停位置、用户发射功率和时隙外队列状态，集中研究一个时隙内全频复用无线调度与硬截止算力准入的耦合关系。

上述边界并非对动态飞行问题的简化替代，而是对研究对象的分层：飞行层给定悬停位置和可用算力，时隙控制层决定实际服务 UAV、子载波、NOMA 角色、卸载比例和 MEC CPU。由此，本文的算法结论仅对应时隙控制层，不能外推为轨迹、功率和电池的联合最优解。

\section{分组 NOMA、用户配对与干扰外部性}
NOMA 允许多个用户复用同一时频资源，其性能依赖信道差异、功率配置、解码顺序和 SIC 可行性~\cite{ding2017survey,ali2016dynamic}。已有子载波分配和交换匹配方法主要针对单小区或预先确定关联关系下的用户调度~\cite{di2017sub,bayat2016matching}。在多 UAV 全频复用上行中，同一子载波不仅包含服务 UAV 内部的 NOMA 叠加信号，还包含其他 UAV 所服务用户在该接收端形成的同频干扰。因此，一个用户改变服务 UAV、子载波或 NOMA 角色后，可能同时改变多个非参与链路的 SINR、SIC 状态和可用速率。

匹配理论将这种“一个匹配改变第三方收益”的现象称为外部性；存在外部性时，双边稳定通常不等价于全网效率~\cite{sasaki1996two,pycia2012stability}。本文不把一般匹配结论直接当作无线算法的性能证明，而是将其转化为可执行的网络级评价：每次候选移动后，按照实际服务 UAV 重新计算全网接收信号、跨 UAV 干扰、SIC 可行性、任务时延和 MEC 负载，并且只接受全网字典序目标严格改善的候选。由此得到的是给定邻域内的确定性局部最优，而不是一般规模上的全局最优保证。

\section{近年多 UAV NOMA-MEC 方法比较}
表~\ref{tab:recent-noma-uav-mec} 按场景、决策变量、目标函数和任务语义比较 2024--2026 年与本文最接近的六项工作。比较的目的不是按期刊层级排列算法，而是判断其原始算法能否在本文模型中保持相同的决策空间和优化目标。

\begin{table}[htbp]
  \centering
  \caption{2024--2026 年近邻多 UAV NOMA-MEC 工作与本文问题口径比较}
  \label{tab:recent-noma-uav-mec}
  \scriptsize
  \renewcommand{\arraystretch}{1.15}
  \begin{tabular}{p{1.8cm}p{3.2cm}p{4.1cm}p{4.6cm}}
    \toprule
    文献 & 场景与目标 & 原方法的核心决策 & 与本文口径的主要差异 \\
    \midrule
    Chen 等~\cite{chen2026qos} & 多 UAV 协作 NOMA-MEC；面向任务优先级与 QoS 的长期系统效用 & UAV 三维轨迹、移动用户关联、卸载比例和资源分配 & 轨迹和长期状态是其策略动作的重要组成；加权效用不同于本文单位价值任务的单时隙硬截止完成基数 \\
    Wu 等~\cite{wu2025multiagentnoma} & 三维动态 UAV-MEC-NOMA；在能耗约束下最小化系统总时延 & 解析卸载、用户聚类、UAV 轨迹和功率分配 & 当前模型不优化轨迹和功率；总时延目标与按时完成任务数目标不可直接互换 \\
    Lin 等~\cite{lin2025drlresource} & 多 UAV 应急 MEC；在时延、能耗和速率约束下提高成功卸载任务数 & 服务 UAV、NOMA 载波、用户功率、动态频谱接入和 UAV 轨迹 & 目标与本文最接近，但其混合动作、长期 MDP 和动态频谱口径不同；冻结核心动作后只能构成适配版而非原算法复现 \\
    Wang 和 Luo~\cite{wang2024energyefficient} & UAV 间 FDMA、UAV 内 NOMA；最小化加权能耗 & 发射功率、计算资源、任务数据分配和 UAV 轨迹 & 跨 UAV 正交频谱消除了本文研究的同频外部性；任务可拆分到多个 UAV，也不符合本文唯一服务 UAV 语义 \\
    Lyu 等~\cite{lyu2024computing} & 海洋物联网 NOMA-UAV 应急通信；最小化完成时延与终端能耗的加权开销 & 发射功率、UAV 算力、任务卸载和载波联盟 & 联盟交换思想可用于构造机制基线，但目标、功率变量和跨接收端干扰语义均不相同 \\
    Liu 等~\cite{liu2025trajectory} & 双 UAV 协作 NOMA-MEC；在能耗约束下最小化系统时延 & UAV 间任务迁移、卸载比例和 UAV 轨迹 & 其“跨 UAV 协作”指 UAV 间任务转移；本文 CrossUAV 指地面用户改变实际服务接收机，时延链和决策语义不同 \\
    \bottomrule
  \end{tabular}
\end{table}

六项工作均至少在可控变量、物理层、任务语义、目标函数或时间尺度中的一个方面与本文不一致。若直接冻结轨迹、功率或动态带宽，原算法的主要动作维度及状态转移将被删除；若保留这些动作，又会使其拥有本文算法没有的决策空间，形成不公平比较。本项目当前尚未取得能够审计相关算法全部训练配置、停止条件并可直接接入本文评价器的作者实现。因此，本文不搬用异构场景中的已发表数值，也不把自行补全的深度强化学习实现标注为“复现 SOTA”。其中，Lin 等提出的 POS\_MADDPG\_NOMA 因目标中包含成功卸载任务数，可作为后续优先申请代码并开展静态适配的候选；在完成预注册适配和独立验证前，仍只作结构化文献比较。

\section{截止时间驱动的卸载与资源准入}
NOMA-MEC 资源分配通常同时包含连续卸载变量和离散任务准入~\cite{ding2019impact,zhu2022multi}。近期研究还从峰值信息年龄边界~\cite{chen2024peakaoi}、异构 MEC 的时延与能耗权衡~\cite{islam2025delta}等角度建立截止或时效约束。这些方法说明任务时限不能仅作为结果统计项，而应进入卸载和算力分配过程。

本文关注单位价值任务在逐 UAV 硬 CPU 预算下的最大按时完成基数。DD-MCAA 先由统一时延模型反推每项任务的最小卸载比例和最小 MEC CPU 需求，再按照最小 CPU 需求与任务编号的稳定顺序选择最长可行前缀。该排序结论依赖任务价值相同、每项任务只由一个 UAV 执行和单时隙 CPU 硬预算；对于加权收益、跨时隙迁移、抢占式队列或多 UAV 分割计算，不能直接沿用该最优性结论。

\section{可比性判据与三层证据定位}
为避免把“算法名称相似”误当作“问题口径相同”，本文只在以下条件同时成立时把文献算法视为直接基线：可控变量一致，均采用相同的跨 UAV 全频复用与硬 SIC 模型，任务均按唯一到达和唯一服务 UAV 计数，主目标均为按时完成任务数，算法获得相同 CSI 与资源预算，并且算法步骤、超参数、停止条件和训练预算可审计。不能满足这些条件的近邻工作保留在表~\ref{tab:recent-noma-uav-mec} 中作研究定位，不进入数值排序。

本文的实验比较据此分为三个层次。

\begin{enumerate}
  \item \textbf{A 层：受控策略。} Fixed-half/Equal 与 Crossing/Equal 仅用于隔离 DD-MCAA 中“时限反推资源需求”和“最大基数准入”的作用，是机制消融和朴素控制组，不代表现有先进方法。
  \item \textbf{B 层：同口径强机制基线。} IA-greedy、per-UAV ILP、exchange-stable 和公平 OFDMA 与本文共享不可变任务/信道轨迹、硬 SIC、逐 UAV CPU 预算和最终评价器，用于检验全网评价、搜索细化、交换稳定机制和接入方式的差异。其中 exchange-stable 只代表联盟博弈/交换匹配类局部稳定机制，并非对 Lyu 等算法的逐式复现。
  \item \textbf{C 层：小规模精确参照。} 对预声明的小规模实例完整枚举“未服务/UAV/子载波”分配，并用同一硬评价器筛除容量、SIC、唯一任务和 CPU 违规。该层用于报告启发式与精确最优的经验差距，但不能证明一般规模的全局最优性。
\end{enumerate}

三层证据分别回答“机制是否有效”“在相同问题口径下是否优于有竞争力的替代机制”和“小规模时距精确最优有多远”。它们不能相互替代：受控策略不能包装为 SOTA，小规模枚举也不能替代一般规模的复杂度或近似比证明。

\section{本文研究定位}
本文不把“多 UAV + NOMA + MEC”的场景首创性作为主要创新，因为近年研究已经覆盖多 UAV 协作、NOMA 接入、轨迹控制和资源分配的多种组合。本文的贡献定位于一个边界清晰的时隙控制问题：在固定悬停、全频复用、硬 SIC、唯一任务计数和有限逐 UAV CPU 条件下，利用全网候选重评价处理跨接收端干扰外部性，并利用截止时间反推的最小资源需求完成最大基数准入。相应创新属于算法机制和可审计优化流程，而场景模型用于限定结论成立的工程条件。

本章的比较同时给出了后文实验解释的边界：只有 A/B/C 三层中的同口径证据用于支持性能主张；异构文献用于说明研究脉络和变量差异，不用于宣称本文全面优于现有最先进方法。
